\section{Measurement and Collapse}


本節では、測定が意識状態、とくにクオリアに対して何を行うのかを整理する。

ここでいう collapse は、量子力学における波動関数の収縮をそのまま意味するものではない。

本稿でいう collapse とは、ライブな境界状態が、報告可能で保存可能なログへ変換されることを指す。

すなわち、測定とは、意識を見つける操作ではなく、意識の開口状態を変形し、低速層へ沈殿させる操作である。

\subsection{Measurement Requires Stabilization}

測定には安定化が必要である。

何かを測定するためには、それが比較可能であり、反復可能であり、記録可能でなければならない。

神経活動を測る場合でも、行動を測る場合でも、主観報告を取る場合でも、対象は何らかの形で安定化される。

しかし、本稿が扱う意識状態は、完全に安定化された対象ではない。

意識は、速度差、結合粘性、内部遅延、意味勾配が一定の範囲に入ったときに開かれる、時間的非閉包のレジームである。

このレジームは、静止した対象ではなく、動的な関係である。

したがって、測定は中立的な観察ではない。

測定は、動いている関係を、記録可能な形へ変換する操作である。

\subsection{Measurement as Plasticity Closure}

測定が行う最も重要な操作は、可塑性ゲートの閉鎖である。

クオリアは、高速層の処理が低速層へ沈殿しつつある境界に現れる。

この境界は、まだ完全にはログ化されていない。

しかし、測定はその境界を報告可能な形に変える。

報告する。

選択肢から選ぶ。

強度を数値化する。

脳活動として記録する。

言語で説明する。

これらはいずれも、ライブな境界状態を低速層へ沈殿させる操作である。

したがって、測定はクオリアをそのまま取り出すのではない。

測定は、クオリアをログ化する。

このとき、クオリアはクオリアとしては消え、代わりに残余が生じる。

\subsection{From Qualia Live to Sedimented Residue}

この変換は、次のように表せる。

\begin{equation}
\text{qualia live}
\rightarrow
\text{measurement}
\rightarrow
\text{sedimented residue}
\end{equation}

ここで qualia live とは、沈殿前線に現れているライブな前景である。

sedimented residue とは、その境界状態が測定によって低速層へ沈殿した後に残るログ、報告、ラベル、数値、神経的要約である。

重要なのは、残余が無意味ではないという点である。

測定後に残るログは、確かに分析できる。

比較もできる。

統計処理もできる。

しかし、それはクオリアそのものではない。

それは、クオリアが通過した後の痕跡である。

したがって、測定可能なものは、クオリアの本体ではなく、クオリアの通過跡である。

\subsection{Why Precision Makes Qualia Disappear}

測定精度を高めるほど、クオリアが消えるように見えることがある。

これは、測定が不完全だからではない。

むしろ、測定が成功しているからである。

精密な測定は、対象をより強く安定化する。

安定化が強まるほど、ライブな境界は低速層へ沈殿しやすくなる。

その結果、測定によって得られるものは、より明確なログ、より明確なラベル、より明確な数値になる。

しかし、その明確さは、クオリアそのものの明確さではない。

それは、クオリアが沈殿した後の残余の明確さである。

このため、測定はクオリアを明らかにする一方で、クオリアをクオリアとしては消してしまう。

ここに、クオリア測定の逆説がある。

\subsection{Measurement and the Consciousness Window}

この測定による変形は、consciousness window の観点からも理解できる。

安定した意識状態は、次のような中間領域に存在する。

\begin{equation}
0 < |\Delta v| < \Delta v_{\max}
\end{equation}

\begin{equation}
\mu_{\min} < \mu < \mu_{\max}
\end{equation}

\begin{equation}
\tau_{\min} < \tau < \tau_{\max}
\end{equation}

\begin{equation}
\mathcal{T}_{\min}
<
\mathcal{T}_{\Phi}
<
\mathcal{T}_{\max}
\end{equation}

測定は、この窓の内部にあるライブな状態を、より低速で安定したログへ移動させる。

つまり、測定は $\Delta v$、$\tau$、$\mu$、$\mathcal{T}_{\Phi}$ の関係を変える。

とくに、報告や記録は、境界状態を低速層へ寄せる。

その結果、測定対象はライブな開口から、沈殿した残余へ変換される。

この変換を非閉包の側から見れば、collapse は速度差の消失として表される。

\begin{equation}
\Delta v \to 0
\quad \Rightarrow \quad
C(t) \to 0
\end{equation}

速度差が消失すると、ログ参照の開きは閉じる。

したがって、collapse は意識の生成ではなく、非閉包状態の終了として解釈できる。

これは量子論そのものの再解釈ではない。

本論文における意識構造の側から見た対応である。

この意味で、測定は外部から意識をそのまま覗く操作ではない。

測定は、意識状態のレジームを変更する操作である。

\subsection{Measurement Is Not Error}

ここで注意すべきことがある。

測定がクオリアを残余へ変換するからといって、測定が無意味になるわけではない。

測定は失敗ではない。

測定は、ライブな境界状態を扱うために必要な変換である。

問題は、測定結果をクオリアそのものと取り違えることである。

測定によって得られる主観報告、反応時間、神経活動、言語記述は、すべて重要である。

しかし、それらは境界現象そのものではなく、境界現象が沈殿した後のログである。

したがって、本稿の立場は反測定主義ではない。

むしろ、測定が何を測っているのかを正確に置き直す立場である。

測定はクオリアを測るのではない。

測定は、クオリアが通過した後に残る沈殿構造を測る。

\subsection{Introspection as Internal Measurement}

この構造は、外部測定だけでなく内省にも当てはまる。

内省とは、意識状態を自分自身で観察し、言語化し、把握しようとする操作である。

しかし、内省もまた測定である。

ある感覚に注意を向ける。

それを言葉にする。

強さを評価する。

原因を考える。

これらの操作は、ライブな境界状態を低速層へ沈殿させる。

したがって、内省によって捉えられるものも、クオリアそのものではなく、クオリアが沈殿した後の構造である。

内省が重要でないという意味ではない。

内省は、低速層におけるログ更新の強力な方法である。

ただし、内省はクオリアを保存するのではなく、クオリアをログ化する。

この点を取り違えると、内省は意識の本体を直接見ているように錯覚される。

\subsection{Experimental Implication}

本稿の枠組みでは、意識研究における実験設計も少し異なる見方を持つ。

測定すべきなのは、クオリアそのものではない。

測定すべきなのは、クオリアがどのような条件で現れ、どのような条件で沈殿し、どのような残余を残すかである。

したがって、実験的に重要なのは、次のような量である。

\begin{enumerate}
\item ライブな経験が報告可能なログへ変換されるまでの遅延。
\item 報告要求の強さによる経験内容の変形。
\item 意味負荷の違いによる残余の変化。
\item 注意操作による $\tau$ や $\mathcal{T}_{\Phi}$ の変化。
\item 報告前後での神経活動パターンの変化。
\item 言語化によって経験がどのように安定化されるか。
\end{enumerate}

これらは、クオリアを直接捕まえる試みではない。

クオリアが沈殿する過程を調べる試みである。

\subsection{Relation to the Measurement Problem}

本節で述べた collapse は、量子測定問題の直接的な解釈ではない。

しかし、構造的な類似はある。

どちらの場合も、測定は対象をそのまま保存する中立的操作ではない。

測定は、状態を安定化し、記録可能な形に変える操作である。

クオリアの場合、この操作はライブな境界状態を沈殿残余へ変換する。

したがって、意識研究における測定問題とは、クオリアがどこに隠れているかという問題ではない。

それは、測定によって境界状態がどのようにログ化されるかという問題である。

\subsection{Summary}

本節では、測定を意識状態の外部観察ではなく、開口状態の変換として整理した。

\begin{itemize}
\item クオリアは沈殿前線に現れるライブな境界状態である。
\item 測定は、その境界状態を低速層へ沈殿させる。
\item 測定後に残るのは、クオリアそのものではなく sedimented residue である。
\item 測定精度が上がるほど、残余は明確になるが、ライブなクオリアは閉じられる。
\item 内省もまた内部測定であり、クオリアをログ化する。
\item 測定は誤りではなく、境界状態を扱うための変換である。
\item 実験的に扱うべきなのは、クオリアそのものではなく、クオリアが沈殿する条件と残余である。
\end{itemize}

したがって、測定とは、意識をそのまま捕まえることではない。

測定とは、時間的非閉包の開口を閉じ、ライブな境界を低速層のログへ変換する操作である。

次節では、この性質が、なぜ意識を特定の場所へ局在化できないのかという問題へつながることを示す。
